The common cubic potential gives the following pullback construction,
including for targets with nonisolated singularities.

\begin{corollary}[Pullback of the common cubic potential]
\label{cor:common-cubic-pullback}
Let \(f\not\equiv0\) be a harmonic cubic on \(\R^m\), and put
\(q=|y|^2\), \(S=|\nabla_yf|^2\), and \(H=f\sqrt{qS}\).
Suppose, on \(S>0\), that
\[
 f\Gamma_y(f,S)\geq0,\qquad H\cL_y(H)\geq0.
\]
For every quadratic Clifford map \(P:\R^N\to\R^m\) with
\(N\geq32\), set \(F=f(P)\), \(Q=q(P)\), and \(T=S(P)\).
Then
\[
 G=R^2F\sqrt{QT}
\]
is a subcalibration potential for \(\{F<0\}\). Its boundary current
is locally mass minimizing and strictly area minimizing, with support
\(\{F=0\}\).
The defining polynomial has degree six and the potential has
homogeneity degree sixteen.
\end{corollary}

\begin{proof}
Put \(H_0=f\sqrt S\). Bochner's identity and \(\Delta_y f=0\) give
\[
 H_0\Delta_yH_0
 =f^2\left(|\Hess f|^2-
       \frac{|\Hess f(\nabla f)|^2}{S}\right)
       +f\Gamma_y(f,S)\geq0.
\]
Since \(H_0\) has degree five, differentiating \(H=\sqrt q\,H_0\)
yields
\[
 H\Delta_yH=qH_0\Delta_yH_0+(m+9)H_0^2\geq0.
\]
Thus the ratio condition holds with \(M=0\). Euler's identity gives
\(\nabla H\ne0\) where \(H\ne0\); at a zero with \(S>0\),
\(\nabla H=\sqrt{qS}\nabla f\ne0\).
The exceptional set \(Z=\{\nabla f=0\}\) is conical and algebraic,
and Euler's identity gives \(Z\subset\{f=0\}\).
Its codimension is at least two by harmonicity and
\cref{lem:square-free-critical-set}.
Now apply \cref{thm:clifford-area-transfer} with \(\ell=6\),
\(a=2\), and \(M=0\): its bound is \(N\geq1+24+4+3=32\).
Harmonicity gives both signs near every target zero and hence the
support and strict-minimality assertions.
\end{proof}

The first hypothesis holds for the radial cubic models, where
\(\Gamma(f,S)\) is a positive multiple of \(qf\), and for the
nonradial models of \cref{thm:nonradial-models}, whose quadratic
minimality weight is positive. The second holds in the minimizing ranges of
\cref{thm:cubic-unified-sextic,thm:nonradial-minimizing}.
Since \(q,f,S\in\cA_q(f)\), this potential is a function of the
smallest target algebra. For an essential target with \(r\) minimal
homogeneous generators and a pullback with \(Q<R^2\) somewhere,
\cref{thm:smallest-clifford-pullback} gives exactly \(r+1\) source
generators. Thus the three-, four-, and five-generator cubic models
yield four-, five-, and six-generator pullback algebras whenever the
two displayed inequalities hold.

For the minimizing nonradial target in \cref{thm:nonradial-minimizing},
the degree-five potential \(H_0=f\sqrt S\) already satisfies
\(H_0\cL(H_0)\ge0\), and the Bochner identity above gives
\(H_0\Delta H_0\ge0\). Thus \cref{thm:clifford-area-transfer}, with
\(\ell=5\), \(a=1\), and \(M=0\), gives the degree-twelve potential
\(G=RF\sqrt T\) for every such Clifford map with \(N\ge28\).

For Cartan targets, a weighted inequality gives the lower-degree
pullback \(G=H(P)\).
